\chapter{Практическое задание 6. Нагрев пластины внутренним источником}

\section{Постановка задачи}

Рассмотрим однородную неограниченную пластину толщины $2R$.
Выберем координату $x$ поперёк пластины, с началом в её средней плоскости:
$-R<x<R$. Из однородности условий вдоль граней температура зависит
только от $x$ и времени $t$.

Начальная температура нулевая. С момента $t=0$ действует постоянный
равномерный объёмный источник тепла плотности $Q$ (энергия на единицу
объёма за единицу времени). Обе грани поддерживаются при нулевой
температуре. Обозначим теплопроводность через $k>0$, объёмную
теплоёмкость через $c>0$ и положим
\[
 a^2=\frac{k}{c},\qquad F_0=\frac Qc.
\]
Тогда
\[
\boxed{\begin{cases}
 u_t=a^2u_{xx}+Q/c,&-R<x<R,\quad t>0,\\
 u(-R,t)=u(R,t)=0,\\
 u(x,0)=0.
\end{cases}}
\]
Если $Q$ уже обозначает нормированную правую часть уравнения,
вместо $Q/c$ следует писать $Q$; здесь используется физическая
объёмная плотность тепловыделения.
Решение чётно по $x$. На половине пластины можно эквивалентно
задать $u_x(0,t)=0$, $u(R,t)=0$.

\section{Стационарное распределение}

Стационарная температура удовлетворяет
\[
 kU''+Q=0,\qquad U(-R)=U(R)=0.
\]
Интегрируя и используя симметрию, получаем
\[
\boxed{U(x)=\frac Q{2k}(R^2-x^2).}
\]
При $Q>0$ максимум достигается в средней плоскости:
\[
 U(0)=\frac{QR^2}{2k}.
\]
Температура конечна: выделяющееся тепло отводится через обе грани.

\section{Решение методом Фурье}

Положим $u=U+v$. Тогда
\[
 v_t=a^2v_{xx},\qquad v(-R,t)=v(R,t)=0,
 \qquad v(x,0)=-U(x).
\]
Для чётного решения собственные функции имеют вид
\[
 X_n(x)=\cos(\eta_nx),\qquad
 \eta_n=\frac{(n+\tfrac12)\pi}{R},\qquad n=0,1,\ldots.
\]
Они удовлетворяют $X_n(\pm R)=0$ и нормировке
\[
 \int_{-R}^RX_n^2(x)\,dx=R.
\]
Нечётные собственные функции не участвуют, поскольку начальные данные
и источник чётны.
Разложим стационарный профиль:
\[
 U(x)=\sum_{n=0}^{\infty}U_n\cos(\eta_nx),\qquad
 U_n=\frac1R\int_{-R}^RU(x)\cos(\eta_nx)\,dx.
\]
Поскольку $\cos(\eta_nR)=0$, $\sin(\eta_nR)=(-1)^n$,
\[
 \int_0^R(R^2-x^2)\cos(\eta_nx)\,dx
 =\frac{2(-1)^n}{\eta_n^3}.
\]
Поэтому
\[
 \boxed{U_n=\frac{2Q(-1)^n}{kR\eta_n^3}
 =\frac{16QR^2}{k\pi^3}\frac{(-1)^n}{(2n+1)^3}.}
\]
Начальная функция $v=-U$ имеет противоположные коэффициенты;
следовательно,
\[
 v(x,t)=-\sum_{n=0}^{\infty}U_n e^{-a^2\eta_n^2t}\cos(\eta_nx).
\]
Получаем две эквивалентные формы полного решения:
\[
\boxed{\begin{aligned}
 u(x,t)={}&\frac Q{2k}(R^2-x^2)\\
 &-\frac{16QR^2}{k\pi^3}\sum_{n=0}^{\infty}
 \frac{(-1)^n}{(2n+1)^3}
 e^{-a^2(2n+1)^2\pi^2t/(4R^2)}
 \cos\frac{(2n+1)\pi x}{2R},
\end{aligned}}
\]
или
\[
\boxed{\begin{aligned}
 u(x,t)=\frac{16QR^2}{k\pi^3}\sum_{n=0}^{\infty}
 \frac{(-1)^n}{(2n+1)^3}
 &\left(1-e^{-a^2(2n+1)^2\pi^2t/(4R^2)}\right)\\
 &\times\cos\frac{(2n+1)\pi x}{2R}.
\end{aligned}}
\]
Если координата отсчитывается от левой грани, $0\le y\le2R$,
нужно заменить $x$ на $y-R$.

\section{Проверка и предельный режим}

На гранях $x=\pm R$ каждое собственное слагаемое равно нулю;
при $t=0$ во второй форме решения множители $1-e^0$ равны нулю.
Уравнение проверяется либо подстановкой первой формы, либо
решением модальных уравнений для постоянного источника.
При $t>0$ ряд допускает необходимое дифференцирование.

При $t\to\infty$ переходная часть затухает и
\[
\boxed{u(x,t)\longrightarrow\frac Q{2k}(R^2-x^2).}
\]
Сходимость равномерная. Её ведущая скорость определяется первой модой:
\[
 u(x,t)-U(x)=-\frac{16QR^2}{k\pi^3}
 e^{-a^2\pi^2t/(4R^2)}\cos\frac{\pi x}{2R}
 +O\left(e^{-9a^2\pi^2t/(4R^2)}\right).
\]
Для $Q>0$ температура неотрицательна и не убывает по времени;
её максимум по толщине достигается в средней плоскости.
Средняя стационарная температура равна
\[
 \overline U=\frac1{2R}\int_{-R}^RU(x)\,dx=\frac{QR^2}{3k}.
\]
Наружные плотности потоков на обеих гранях одинаковы:
\[
 j_{\mathrm{out},+}=-kU'(R)=QR,\qquad
 j_{\mathrm{out},-}=kU'(-R)=QR.
\]
Их сумма равна $2QR$, то есть тепловыделению в пластине на единицу
площади её грани. Это проверяет стационарный энергетический баланс.
При $Q=0$ решение тождественно нулевое.
